\subsection{Real-Robot Refinement on Stair-Based Manipulation}
\label{sec:exp-real-robot}

We evaluate a pretrained flow-matching policy on the OpenArm left arm
and parallel gripper, using chest and wrist RGB observations together with
joint and gripper state. The tasks are to place Jigglypuff into an adjacent
case and to stack a red cube on a purple cube
(Fig.~\ref{fig:real-robot}). Both use a staircase: the stepped surfaces
introduce different approach heights and nearby edges, making grasping and
placement sensitive to alignment. Stacking additionally requires precise
positioning and release without disturbing the supporting cube. This task
design exposes a practical refinement problem: useful reaching, grasping,
and transport can still end in an unsuccessful placement or release.

We compare the pretrained \textbf{Base}, supervised \textbf{BC} refinement,
\textbf{DICE-RL} with its original actor objective, and \textbf{CAST} with
constrained action targets. Policy pretraining uses 40 placement
demonstrations and approximately 60 stacking demonstrations. For each
task, BC refinement uses 30 additional expert demonstrations. The two RL
methods share 40 recorded robot rollouts per task, retaining both
successful and failed attempts. Refinement then trains offline on the
saved data. BC fine-tunes the base for 40 additional epochs. Each RL run learns a residual with the base and
encoder frozen, using 6,000 critic updates. DICE-RL and CAST share the base, replay data and
minibatch schedule, critic-training procedure, and update budget. BC
therefore differs from RL in both supervision and trainable parameters.

After refinement, we evaluate each method in 20 autonomous attempts per
task. Placement poses are randomized on the stairs, and success requires
the toy to remain fully inside the case. Stacking succeeds when the red
cube remains on the purple cube after release and gripper withdrawal.
Detailed protocols are in the experimental supplement.

CAST raises placement success from 55\% to 95\% and stacking success from
35\% to 70\%, gains of 40 and 35 percentage points over the respective
bases. These gains support H1 under offline refinement and physical
deployment. Fig.~\ref{fig:real-robot} reports all four methods and their
successful attempt counts; CAST achieves the highest observed success on
both tasks.

The main observed stacking failure is failure to release the gripper after
useful reaching, grasping, and alignment. Refinement should retain those
useful actions while correcting the release. BC receives explicit release
targets from expert demonstrations; CAST moderates critic-guided changes
through restoration and clipping. The observed ordering is consistent
with preserving useful actions while correcting task-ending mistakes.
